\begin{figure}[H]
  \centering
  \resizebox{0.95\textwidth}{!}{%
    \begin{tikzpicture}[                                                                                                                  
        font=\footnotesize,
        bloque/.style={draw=diagNodoBorde, fill=diagNodo, align=left,
            inner sep=7pt, rounded corners=2pt}
      ]
      % Caso real: bloque de salidas del Grado en Ingenieria Informatica
      % (1.137 caracteres en data/grados.json, rastreo del 05/08/2026).
      \node[bloque, text width=10.5cm] (txt) {%
        \textbf{Campo \texttt{texto}} \hfill \textit{un único bloque, sin estructura interna}\\[4pt]
        «El graduado en Ingeniería Informática presenta multitud de salidas
        profesionales. Es una de las titulaciones que tiene más demanda en el
        mercado laboral, con una tasa de paro inexistente y una oferta de
        trabajo que crece cada día. El graduado en Ingeniería Informática
        desempeña sus funciones en todos los sectores de la actividad
        económica, en empresas, consultoras informáticas, la
        administración\dots»};
      \node[bloque, text width=5.2cm, right=8mm of txt.north east,
        anchor=north west] (met) {%
        \textbf{Metadatos}\\[4pt]
        \texttt{tipo}: salidas\\
        \texttt{grado}: Grado en Ingeniería Informática\\[4pt]
        No lleva código: la unidad es la \textbf{titulación entera}, no una
        asignatura.};
      \node[bloque, fill=diagGrupo, draw=diagGrupoBorde, text width=16.6cm,
        below=6mm of txt.south west, anchor=north west] (nota) {%
        La colección tiene \textbf{8 bloques para 12 titulaciones}: los siete
        grados simples y un único doble grado. Los otros cuatro dobles no
        publican salidas propias.};
    \end{tikzpicture}%
  }
  \caption[Estructura de un bloque de salidas profesionales]{\textcolor{borrador2}{Estructura
      de un bloque de salidas profesionales, sobre el caso real del Grado en Ingeniería
      Informática. Fuente: Elaboración propia.}}
  \label{fig:estructura_salidas}
\end{figure}
